% ============================================================
% 分片：s01-intro  对应原文 PDF 第 1--3 页
% 负责章节：§1 Introductory remarks（第 1 页为标题页，标题/作者/摘要由 main.tex 承担）
% 译者：translator-01   审校状态：初稿
% ============================================================

% ==== 原文 PDF 第 1 页 ====
% [说明: 第 1 页为标题页（标题、作者、单位、摘要与会议信息），已由 main.tex 承担，此处不重复。]

% ==== 原文 PDF 第 2 页 ====
\section{引言}

\begin{keytakeaway}
\textbf{核心物理图像：}对偶是保持可观测量（底层物理）不变的等价关系；它通常在任何一个弱耦合描述中都不显现，
因而在微扰论里不易察觉。许多弦对偶把一个理论的强耦合极限与另一个（一般不同的）理论的弱耦合描述等同起来。

\textbf{数学工具：}把对偶类比为傅里叶变换；在超对称情形下路径积分可局域化，受保护量不依赖耦合常数，
于是可在任一框架中计算；其数学化身由 Fourier--Mukai 变换、Koszul 对偶等傅里叶变换的推广来刻画。

\textbf{主要结论：}本讲义以扭变、上同调与高阶代数为线索，研究弦论与对偶的数学侧面；镜像对称是贯穿全篇的原型对偶。
\end{keytakeaway}

本讲义基于 TASI 2021 上所作的报告，该系列报告旨在概述弦论中的基本对偶，以及用于检验和探索这些对偶的数学技术。
具体而言，本讲义着重于该系列讲座中最后两讲的内容，这两讲意在凸显某些在当代物理数学与超对称弦论研究中无处不在的数学结构与框架。
我们所介绍的技术已被证明对于形式化与概念化这些理论都不可或缺，并在检验各种拟议中的对偶（putative dualities）时得到了广泛使用。
此外，它们本身就引出了大量独立的学术探究。
本综述力求自足，同时对所选主题给出简洁的教学式导引；因此，主题（以及文献）的选取非常有限，而且必然是不完整的。
该系列讲座的前两讲在本讲义中不予回顾，它们对弦论作了闪电式的入门介绍，并概述了各种背景上的基本超对称弦对偶。
我们略去它们，是因为这些主题已在别处（例如 [1]）以及各种弦论基础教材中有详尽讨论。

在深入正题之前，我们先依照 [2] 勾勒若干一般性的哲学。
对偶是等价：即一个（通常极为复杂的）变量变换，它使底层的物理，即可观测量，保持不变。
对偶之所以不同于较为平凡的基变换、之所以非平凡，是因为它们通常在任何弱耦合描述中都不显现。
毋宁说，它们是理论精确的「对称性」。
\footnote{在本讲义中，我们关注精确对偶，而非红外对偶。}
这里之所以给「对称性」加上引号，是因为对偶除了包含场变量的变换之外，往往还涉及背景参数（例如耦合常数）的非平凡改变。
\footnote{然而要注意，在弦论中这些参数是由模场的真空期望值设定的，这一点应与场论中的情形相对照。}
在许多弦对偶中，一个弦论的强耦合极限通过对偶与另一个（一般不同的）弦论的弱耦合描述相联系（即等价）。

% [Tutor Note: 严格说，$e^{iS/(g^2\hbar)}$ 是路径积分的权重（作用量的指数化），作者把「作用量的权重」略称为经典作用量；
% 关键之处在于它只依赖组合 $g^2\hbar$，因此「$g$ 小」可与「$\hbar$ 小」等同起来，这正是下文弱–强对偶与傅里叶变换类比的出发点。]
粗略地说，可以把这类对偶看成类似于傅里叶变换的东西。回忆一下，可以把弱耦合的（比如说）弦论的经典作用量写成
\[
e^{iS/(g^2\hbar)}
\]
其中我们已经对场作了重新标度，使耦合常数作为整体因子 $1/g^2$ 出现在作用量之前。
由于耦合常数此后只以组合 $g^2\hbar$ 的形式出现，我们看到小的 $g^2$ 等价于小的 $\hbar$，而大的 $g^2$ 等价于大的 $\hbar$。
一个强–弱对偶把耦合常数变换为 $g\sim 1/g$，于是当我们把 $g$ 从弱耦合调到强耦合时，
就可以把理论原来的描述换成对偶理论的描述，而对偶理论的耦合 $g'\sim 1/g$ 正变得越来越弱。
等价地，也可以把对偶理解为：把量子涨落变得越来越大的描述，换成量子涨落越来越被压制的描述。
这类似于傅里叶变换：在位置空间中高度局域化的量在动量空间中变得弥散，反之亦然。
然而，在强–弱对偶中，变量之间的映射涉及无限维场空间上的高度非线性变换。
一般说来，我们并不知道如何把弱耦合描述中良好的%
% ==== 原文 PDF 第 3 页 ====
% [跨页接续: 原文第2页/3页，该句自上一页续接]
场变量映射到对偶的强耦合一侧的对应物。
这一陈述在场论中已经成立，而在弦论中情形可能更为复杂。

% [Tutor Note: 傅里叶变换的类比只能当作启发：$x$ 与 $p$ 之间是线性（幺正的）变换，而弱–强对偶是无限维场空间上的非线性映射。
% 后文提到的 Fourier--Mukai 变换与 Koszul 对偶，才是这一类比在数学上的精确对应物。]
% [Tutor Note: 「精确对偶」指在全部能标上成立的等价；「红外对偶」只断言低能极限下两理论等价，通常由 RG 流的不动点之间的对应给出，是较弱的陈述。]
% [Tutor Note: 「受保护量」指在耦合常数或背景参数的连续形变下保持不变的量，典型例子是 BPS 态计数（如 Witten 指标）与拓扑不变量。
% 其不变性来自超对称 Ward 恒等式：形变只产生 $Q$ 恰当项，而 $Q$ 恰当的插入在关联函数中为零（见式 (2.1) 的论证）。]
我们将着重讨论保持超对称的理论。
在这些情形下（但并非只在这些情形下），对偶的某些侧面可以再度与数学上良定义的量联系起来。
例如，我们将会看到，路径积分可以局域化到有限维积分，而受保护量可以在耦合常数的形变下（更一般地，在感兴趣的相应对偶下变换的那些参数的形变下）保持不变。
受保护量可以在任何对偶框架中计算，尽管它们在不同框架中的描述可能大相径庭；它们为非微扰物理提供了有用的探针。
人们所提出的这种对偶不变性还会导致高度非平凡的（猜想性的）数学等价。
我们将说明如何借助\emph{扭变}分离出这类受保护量的物理，并讨论支配其算子乘积的若干代数结构。
事实上，当着眼于这类受保护量的数学化身时，某些弦对偶的若干侧面确实可以看作恰好归结为傅里叶变换。
更确切地说，傅里叶变换的各种自然数学推广——例如 Fourier--Mukai 变换，乃至 Koszul 对偶\footnote{关于面向非专业读者的这一视角的阐述，见 [3]。}——无处不在。
尽管在本讲义余下部分我们不会依赖这一视角，但我们建议读者在踏上（超对称）对偶追寻之旅时把这一类比记在心中。

% [Tutor Note: 路径积分「局域化到有限维积分」通常由 Duistermaat--Heckman 定理或 Atiyah--Bott 局域化实现：
% 当作用量是某个超荷的恰当项时，驻相近似是精确的，积分只由 $Q$ 的不动点（或临界点）贡献。
% 对超对称理论而言，这正是受保护量与耦合无关的机制。]
% [背景补充: 「ur-」是德语前缀，意为「原初的、原型的」；「ur-example」即「原型示例」，此处指最能代表主题的范例。]
我们在弦对偶方面的原型（ur-example）——它将在本讲义中始终扮演主角——是镜像对称。
镜像对称是一种微扰弦对偶；它并不是弦耦合常数 $g_s$ 意义上的强–弱对偶。
% [Tutor Note: 半径为 $R$ 的圆上，T 对偶把 $R$ 换成 $\alpha'/R$，于是世界面圈展开参数 $\alpha'/R^2$ 变为 $R^2/\alpha'$；
% 这正是把「大半径」与「小半径」互换，故在 $\alpha'$ 展开的意义上表现为强–弱对偶，而 $g_s$ 在两个框架中都可以很小。]
一个值得记住的有用类比是半径为 $R$ 的圆上的 T 对偶：它在弦论的意义上，即在世界面圈展开参数 $\alpha'/R^2$ 的意义上表现为「强–弱对偶」\footnote{确实，尽管我们不会在本讲义中进一步讨论，镜像对称可以猜想性地理解为 Calabi--Yau 的纤维上 T 对偶的复合 [4]，这里把 Calabi--Yau 视为特殊拉格朗日纤维化。我们还要指出，当 D 膜由凝聚层建模时（如我们在 2.4.2 节所回顾的），T 对偶对 D 膜的作用可以形式化为 Fourier--Mukai 变换。}，
但不在弦耦合 $g_s$ 的意义上——在两个对偶框架中 $g_s$ 都可以保持很小。
其他对偶在作用于 $g_s$ 的方式上将是真正的强–弱对偶，因而是非微扰对偶。
尽管其中某些对偶也许不如镜像对称这样的微扰对偶那样为人所理解\footnote{不过值得强调，弦论可能模糊这些情形之间表面上清晰的区分：例如，基本弦中的非微扰 S 对偶，对于对偶的孤子弦而言就是微扰的 T 对偶 [5]。}，
但这里所介绍的数学概念与技术对这类情形有广泛的适用性。

在本讲义的余下部分，我们将强调\emph{上同调}在研究弦论与对偶的数学侧面时的自然作用，其语境是超对称理论的扭变。
随后我们将讨论\emph{同伦代数}（homotopical algebras），也称为\emph{高阶代数}（higher algebras），它们为扭变理论的物理提供了更为精细的视角。
